\chapter{2014年江苏卷}

\section{填空题}

\begin{problem}
已知集合 \(A = \{-2, -1, 3, 4\}\)，\(B = \{-1, 2, 3\}\)，则 \(A \cap B =\fillinblank{}\).
\end{problem}

\begin{answer}
\(\{-1,3\}\)
\end{answer}

\begin{solution}
交集中的元素必须同时属于集合 \(A\) 和集合 \(B\). 集合 \(A\) 与集合 \(B\) 的公共元素只有 \(-1\) 和 \(3\)，所以
\[
A\cap B=\{-1,3\}.
\]
\end{solution}

\begin{problem}
已知复数 \(z = (5 + 2\i)^2\)（\(\i\) 为虚数单位），则 \(z\) 的实部为\fillinblank{}.
\end{problem}

\begin{answer}
21
\end{answer}

\begin{solution}
利用 \(\i^2=-1\)，得
\[
z=(5+2\i)^2=25+20\i+4\i^2=25+20\i-4=21+20\i.
\]
因此 \(z\) 的实部为 \(21\).
\end{solution}

\begin{problem}
如图是一个算法流程图，则输出的 \(n\) 的值是\fillinblank{}.

\bitmapfigure[width=8.2cm]{img/2014/jiangsu/q3_fig1.png}
\end{problem}

\begin{answer}
5
\end{answer}

\begin{solution}
流程开始时 \(n=0\)，随后每循环一次先令 \(n\leftarrow n+1\)，再判断 \(2^n>20\). 当 \(n=1,2,3,4\) 时，\(2^n\leqslant20\)，继续循环；当 \(n=5\) 时，\(2^5=32>20\)，满足判断条件并输出 \(n\). 所以输出值为 \(5\).
\end{solution}

\begin{problem}
从 \(1\)，\(2\)，\(3\)，\(6\) 这 \(4\) 个数中一次随机地取 \(2\) 个数，则所取 \(2\) 个数的乘积为 \(6\) 的概率是\fillinblank{}.
\end{problem}

\begin{answer}
\(\frac{1}{3}\)
\end{answer}

\begin{solution}
从 \(1\)，\(2\)，\(3\)，\(6\) 中一次取出 \(2\) 个数，基本事件总数为 \(\mathrm C_4^2=6\). 乘积为 \(6\) 的取法只有 \(\{1,6\}\) 和 \(\{2,3\}\)，共有 \(2\) 种，因此所求概率为
\[
\frac{2}{\mathrm C_4^2}=\frac{2}{6}=\frac13.
\]
\end{solution}

\begin{problem}
已知函数 \( y = \cos x \) 与 \( y = \sin (2x + \varphi) \) （\( 0 \leqslant \varphi < \pi \)），它们的图象有一个横坐标为 \( \frac{\pi}{3} \) 的交点，则 \( \varphi \) 的值是\fillinblank{}.
\end{problem}

\begin{answer}
\(\frac{\pi}{6}\)
\end{answer}

\begin{solution}
在 \(x=\frac\pi3\) 处两函数值相等，所以
\[
\sin\left(\frac{2\pi}{3}+\varphi\right)=\cos\frac\pi3=\frac12.
\]
由于 \(0\leqslant\varphi<\pi\)，故
\(\frac{2\pi}{3}\leqslant\frac{2\pi}{3}+\varphi<\frac{5\pi}{3}\). 在这个区间内满足正弦值为 \(\frac12\) 的只有 \(\frac{5\pi}{6}\)，从而
\[
\frac{2\pi}{3}+\varphi=\frac{5\pi}{6}.
\]
从而 \(\varphi=\frac\pi6\).
\end{solution}

\begin{problem}
为了了解一片经济林的生长情况，随机抽测了其中 \(60\) 株树木的底部周长（单位：\(\mathrm{cm}\)），所得数据均在区间 \([80,130]\) 上，其频率分布直方图如图所示，则在抽测的 \(60\) 株树木中，有\fillinblank{} 株树木的底部周长小于 \(100 \mathrm{~cm}\).

\bitmapfigure{img/2014/jiangsu/q6_fig1.png}
\end{problem}

\begin{answer}
24
\end{answer}

\begin{solution}
由直方图可知，底部周长在 \([80,90)\) 和 \([90,100)\) 内的频率分别为 \(10\times0.015=0.15\)，\(10\times0.025=0.25\).
因此底部周长小于 \(100\mathrm{~cm}\) 的频率为 \(0.15+0.25=0.40\)，对应的树木数为
\[
60\times0.40=24.
\]
\end{solution}

\begin{problem}
在各项均为正数的等比数列 \(\{a_{n}\}\) 中，若 \(a_{2} = 1\)，\(a_{8} = a_{6} + 2a_{4}\)，则 \(a_{6}\) 的值是\fillinblank{}.
\end{problem}

\begin{answer}
4
\end{answer}

\begin{solution}
设等比数列的公比为 \(q>0\). 由 \(a_2=1\)，可得
\(a_4=q^2\)，\(a_6=q^4\)，\(a_8=q^6\).
代入已知等式，得
\[
q^6=q^4+2q^2.
\]
令 \(u=q^2>0\)，则
\[
u^3-u^2-2u=u(u^2-u-2)=0.
\]
因为 \(u>0\)，所以 \(u^2-u-2=0\)，即 \((u-2)(u+1)=0\)，从而 \(u=2\). 因此
\[
a_6=q^4=(q^2)^2=4.
\]
\end{solution}

\begin{problem}
设甲，乙两个圆柱的底面积分别为 \(S_{1}\)，\(S_{2}\)，体积分别为 \(V_{1}\)，\(V_{2}\)，若它们的侧面积相等，且 \(\frac{S_{1}}{S_{2}} = \frac{9}{4}\)，则 \(\frac{V_{1}}{V_{2}}\) 的值是\fillinblank{}.
\end{problem}

\begin{answer}
\(\frac{3}{2}\)
\end{answer}

\begin{solution}
设两个圆柱的底面半径和高分别为 \(r_1,h_1\) 以及 \(r_2,h_2\). 由侧面积相等，得 \(2\pi r_1h_1=2\pi r_2h_2\)，从而 \(\frac{h_1}{h_2}=\frac{r_2}{r_1}\).
又因为 \(S_i=\pi r_i^2\)，所以
\(\frac{r_1}{r_2}=\sqrt{\frac{S_1}{S_2}}=\frac32\)，\(\frac{h_1}{h_2}=\frac23\).
因此
\[
\frac{V_1}{V_2}=\frac{S_1h_1}{S_2h_2}=\frac94\times\frac23=\frac32.
\]
\end{solution}

\begin{problem}
在平面直角坐标系 \(xOy\) 中，直线 \(x + 2y - 3 = 0\) 被圆 \((x - 2)^2 + (y + 1)^2 = 4\) 截得的弦长为\fillinblank{}.
\end{problem}

\begin{answer}
\(\frac{2\sqrt{55}}{5}\)
\end{answer}

\begin{solution}
圆心为 \(O_1=(2,-1)\)，半径为 \(R=2\). 圆心到直线 \(x+2y-3=0\) 的距离为
\[
d=\frac{|2+2(-1)-3|}{\sqrt{1^2+2^2}}=\frac3{\sqrt5}.
\]
弦长的一半为
\[
\sqrt{R^2-d^2}=\sqrt{4-\frac95}=\sqrt{\frac{11}{5}}.
\]
所以弦长为
\[
2\sqrt{\frac{11}{5}}=\frac{2\sqrt{55}}5.
\]
\end{solution}

\begin{problem}
已知函数 \( f(x) = x^2 + mx - 1 \)，若对于任意 \( x \in [m, m + 1] \)，都有 \( f(x) < 0 \) 成立，则实数 \( m \) 的取值范围是\fillinblank{}.
\end{problem}

\begin{answer}
\(\left(-\frac{\sqrt{2}}{2},0\right)\)
\end{answer}

\begin{solution}
函数 \(f(x)=x^2+mx-1\) 的图象开口向上，因此在闭区间 \([m,m+1]\) 上的最大值取在端点. 要使任意 \(x\in[m,m+1]\) 都有 \(f(x)<0\)，充要条件是 \(f(m)<0\) 且 \(f(m+1)<0\).
分别计算得
\[
f(m)=2m^2-1<0,
\]
所以 \(-\frac{\sqrt2}{2}<m<\frac{\sqrt2}{2}\)；又
\[
f(m+1)=2m^2+3m=m(2m+3)<0,
\]
所以 \(-\frac32<m<0\). 两个范围取交集，得到
\[
m\in\left(-\frac{\sqrt2}{2},0\right).
\]
\end{solution}

\begin{problem}
在平面直角坐标系 \(xOy\) 中，若曲线 \(y = ax^2 + \frac{b}{x}\)（\(a\)，\(b\) 为常数）过点 \(P(2, -5)\)，且该曲线在点 \(P\) 处的切线与直线 \(7x + 2y + 3 = 0\) 平行，则 \(a + b\) 的值是\fillinblank{}.
\end{problem}

\begin{answer}
\(-3\)
\end{answer}

\begin{solution}
由曲线经过 \(P(2,-5)\)，得
\[
4a+\frac b2=-5.
\]
函数 \(y=ax^2+\frac bx\) 的导数为 \(y'=2ax-\frac b{x^2}\)，所以在 \(P\) 处的切线斜率为 \(4a-\frac b4\). 直线 \(7x+2y+3=0\) 的斜率为 \(-\frac72\)，由平行条件得
\[
4a-\frac b4=-\frac72.
\]
将两式分别化为 \(8a+b=-10\) 和 \(16a-b=-14\).
相加得 \(24a=-24\)，即 \(a=-1\)，再得 \(b=-2\). 因而
\[
a+b=-3.
\]
\end{solution}

\begin{problem}
如图，在平行四边形 \(ABCD\) 中，已知 \(AB = 8\)，\(AD = 5\)，\(\overrightarrow{CP} = 3\overrightarrow{PD}\)，\(\overrightarrow{AP} \cdot \overrightarrow{BP} = 2\)，则 \(\overrightarrow{AB} \cdot \overrightarrow{AD}\) 的值是\fillinblank{}.

\bitmapfigure{img/2014/jiangsu/q12_fig1.png}
\end{problem}

\begin{answer}
22
\end{answer}

\begin{solution}
设 \(\overrightarrow{AB}=\bs u\)，\(\overrightarrow{AD}=\bs v\). 取 \(A\) 为原点，则 \(B=\bs u\)，\(D=\bs v\)，\(C=\bs u+\bs v\). 由 \(\overrightarrow{CP}=3\overrightarrow{PD}\)，得
\[
P=\frac{C+3D}{4}=\frac14\bs u+\bs v.
\]
所以 \(\overrightarrow{AP}=\frac14\bs u+\bs v\)，\(\overrightarrow{BP}=-\frac34\bs u+\bs v\).
由 \(AB=8\)，\(AD=5\) 和题设数量积条件，得
\[
2=\overrightarrow{AP}\cdot\overrightarrow{BP}
=-\frac3{16}|\bs u|^2-\frac12\bs u\cdot\bs v+|\bs v|^2
=-12-\frac12\bs u\cdot\bs v+25.
\]
因此 \(\bs u\cdot\bs v=22\)，即
\[
\overrightarrow{AB}\cdot\overrightarrow{AD}=22.
\]
\end{solution}

\begin{problem}
已知 \( f(x) \) 是定义在 \( \R \) 上且周期为 3 的函数，当 \( x \in [0,3) \) 时，\( f(x) = \left|x^2 - 2x + \frac{1}{2}\right| \)，若函数 \( y = f(x) - a \) 在区间 \([-3,4]\) 上有 10 个零点（互不相同），则实数 \( a \) 的取值范围是\fillinblank{}.
\end{problem}

\begin{answer}
\(\left(0,\frac{1}{2}\right)\)
\end{answer}

\begin{solution}
在 \([0,3)\) 上令
\[
g(x)=\left|(x-1)^2-\frac12\right|.
\]
由于 \(f\) 的周期为 \(3\)，方程 \(f(x)=a\) 的根可由 \(g(x)=a\) 的根平移得到. 对 \(0<a<\frac12\)，方程 \(g(x)=a\) 等价于
\[
(x-1)^2=\frac12+a\quad\text{或}\quad (x-1)^2=\frac12-a.
\]
这给出 \([0,3)\) 内的四个根
\[
x=1\pm\sqrt{\frac12+a}\quad\text{或}\quad x=1\pm\sqrt{\frac12-a}.
\]
其中两个根属于 \((0,1)\)，另外两个根属于 \((1,3)\). 在区间 \([-3,4]\) 中，属于 \([0,1]\) 的每个根有三个周期位置，属于 \((1,3)\) 的每个根有两个周期位置，所以共有
\[
2\times3+2\times2=10
\]
个互不相同的根.

下面说明这是唯一可能的范围. 当 \(a=0\) 时，\(g(x)=0\) 在 \([0,3)\) 内有两个根，按上述周期位置计数得到 \(3+2=5\) 个；当 \(a=\frac12\) 时，根为 \(0,1,2\)，计数为 \(3+3+2=8\) 个. 当 \(\frac12<a<\frac72\) 时，\([0,3)\) 内只有一个根 \(1+\sqrt{\frac12+a}\in(2,3)\)，在 \([-3,4]\) 内只有两个对应根；当 \(a<0\) 或 \(a\geqslant\frac72\) 时没有达到 \(10\) 个根的可能. 因此
\[
a\in\left(0,\frac12\right).
\]
\end{solution}

\begin{problem}
若 \(\triangle ABC\) 的内角满足 \(\sin A + \sqrt{2}\sin B = 2\sin C\)，则 \(\cos C\) 的最小值是\fillinblank{}.
\end{problem}

\begin{answer}
\(\frac{\sqrt{6}-\sqrt{2}}{4}\)
\end{answer}

\begin{solution}
设三角形 \(ABC\) 的边 \(a,b,c\) 分别对应角 \(A,B,C\). 由正弦定理，题设条件等价于
\[
a+\sqrt2b=2c.
\]
令 \(t=\frac ab>0\)，则
\[
\frac cb=\frac{t+\sqrt2}{2}.
\]
三角形不等式给出
\[
\frac{2-\sqrt2}{3}<t<2+\sqrt2;
\]
例如左端来自 \(b-a<c\)（在 \(t<1\) 时），右端来自 \(a-b<c\)（在 \(t>1\) 时），其余不等式自动满足.

由余弦定理，
\[
\cos C=\frac{a^2+b^2-c^2}{2ab}
=\frac{3t}{8}-\frac{\sqrt2}{4}+\frac1{4t}.
\]
记右端为 \(F(t)\)，则
\[
F'(t)=\frac38-\frac1{4t^2}.
\]
因此 \(F(t)\) 在 \(t=\sqrt{\frac23}\) 处取得最小值；该点位于上述取值范围内. 代入得
\[
\cos C_{\min}=\frac38\sqrt{\frac23}-\frac{\sqrt2}{4}+\frac1{4\sqrt{\frac23}}
=\frac{\sqrt6-\sqrt2}{4}.
\]
\end{solution}

\section{解答题}

\begin{problem}
已知 \(\alpha \in \left(\frac{\pi}{2}, \pi\right)\)，\(\sin \alpha = \frac{\sqrt{5}}{5}\).
\begin{enumerate}
\item 求 \( \sin\left(\frac{\pi}{4}+\alpha\right) \) 的值；
\item 求 \( \cos\left(\frac{5}{6}\pi - 2\alpha\right) \) 的值.
\end{enumerate}
\end{problem}

\begin{solution}
因为 \(\alpha\in\left(\frac{\pi}{2},\pi\right)\)，所以 \(\cos\alpha<0\)，从而
\[
\cos\alpha=-\sqrt{1-\sin^2\alpha}=-\sqrt{1-\frac15}=-\frac{2\sqrt5}{5}.
\]

\begin{enumerate}
\item
\[
\sin\left(\frac{\pi}{4}+\alpha\right)=\sin\frac{\pi}{4}\cos\alpha+\cos\frac{\pi}{4}\sin\alpha=\frac{\sqrt2}{2}\left(-\frac{2\sqrt5}{5}+\frac{\sqrt5}{5}\right)=-\frac{\sqrt{10}}{10}.
\]

\item 由
\[
\cos2\alpha=\cos^2\alpha-\sin^2\alpha=\frac45-\frac15=\frac35
\]
和
\[
\sin2\alpha=2\sin\alpha\cos\alpha=-\frac45
\]
得
\[
\cos\left(\frac{5\pi}{6}-2\alpha\right)=\cos\frac{5\pi}{6}\cos2\alpha+\sin\frac{5\pi}{6}\sin2\alpha=-\frac{\sqrt3}{2}\cdot\frac35+\frac12\cdot\left(-\frac45\right)=-\frac{3\sqrt3+4}{10}.
\]
\end{enumerate}
\end{solution}

\begin{problem}
如图，在三棱锥 \(P-ABC\) 中，\(D\)，\(E\)，\(F\) 分别为棱 \(PC\)，\(AC\)，\(AB\) 的中点. 已知 \(PA \perp AC\)，\(PA = 6\)，\(BC = 8\)，\(DF = 5\).

\begin{enumerate}
\item 求证：直线 \(PA \parallel\) 平面 \(DEF\)；
\item 平面 \(BDE \perp\) 平面 \(ABC\).
\end{enumerate}

\bitmapfigure{img/2014/jiangsu/q16_fig1.png}
\end{problem}

\begin{solution}
\begin{enumerate}
\item 在三角形 \(PAC\) 中，\(D\)，\(E\) 分别为 \(PC\)，\(AC\) 的中点，所以由中位线定理得 \(DE\parallel PA\). 若 \(A\in\) 平面 \(DEF\)，则因 \(A,E,C\) 共线且 \(A,F,B\) 共线，平面 \(DEF\) 同时含有直线 \(AC\) 和 \(AB\)，从而平面 \(DEF\) 与平面 \(ABC\) 重合. 于是 \(D\in\) 平面 \(ABC\)，而 \(C,D,P\) 共线，便有 \(P\in\) 平面 \(ABC\)，这与 \(P,A,B,C\) 不共面矛盾. 因此 \(A\notin\) 平面 \(DEF\)，直线 \(PA\) 不在平面 \(DEF\) 内. 又因 \(DE\parallel PA\)，所以 \(PA\parallel\) 平面 \(DEF\).

\item 设 \(\bs{p}=\overrightarrow{AP}\)，\(\bs{b}=\overrightarrow{AB}\)，\(\bs{c}=\overrightarrow{AC}\). 由题设 \(\lvert\bs{p}\rvert=6\)，\(\lvert\bs{c}-\bs{b}\rvert=BC=8\)，且 \(\bs{p}\cdot\bs{c}=0\).
由于 \(D\)，\(F\) 分别为 \(PC\)，\(AB\) 的中点，
\[
\overrightarrow{DF}=\frac12\left(\bs{b}-\bs{p}-\bs{c}\right).
\]
由 \(DF=5\) 得 \(\lvert\bs{p}+\bs{c}-\bs{b}\rvert=10\). 两边平方，得
\[
100=\lvert\bs{p}\rvert^2+\lvert\bs{c}-\bs{b}\rvert^2+2\bs{p}\cdot(\bs{c}-\bs{b})=36+64+2\bs{p}\cdot(\bs{c}-\bs{b}),
\]
所以 \(\bs{p}\cdot(\bs{c}-\bs{b})=0\). 再结合 \(\bs{p}\cdot\bs{c}=0\)，可得 \(\bs{p}\cdot\bs{b}=0\)，即 \(PA\perp AB\). 因为 \(PA\perp AB\)，\(PA\perp AC\)，所以 \(PA\perp\) 平面 \(ABC\).

又因为 \(D\)，\(E\) 分别为 \(PC\)，\(AC\) 的中点，所以 \(DE\parallel PA\). 直线 \(DE\) 在平面 \(BDE\) 内，故平面 \(BDE\) 含有一条垂直于平面 \(ABC\) 的直线，从而平面 \(BDE\perp\) 平面 \(ABC\).
\end{enumerate}
\end{solution}

\begin{problem}
如图，在平面直角坐标系 \(xOy\) 中，\(F_{1}\)，\(F_{2}\) 分别是椭圆 \(\frac{x^2}{a^2} + \frac{y^2}{b^2} = 1\)（\(a > b > 0\)）的左，右焦点. 顶点 \(B\) 的坐标为 \((0, b)\). 连接 \(BF_{2}\) 并延长交椭圆于点 \(A\)，过点 \(A\) 作 \(x\) 轴的垂线交椭圆于另一点 \(C\)，连接 \(F_{1}C\).

\begin{enumerate}
\item 若点 \(C\) 的坐标为 \(\left(\frac{4}{3}, \frac{1}{3}\right)\)，且 \(B F_{2} = \sqrt{2}\)，求椭圆的方程；
\item 若 \(F_{1}C \perp AB\)，求椭圆离心率 \(e\) 的值.
\end{enumerate}

\FigureLayoutDeclare{img/2014/jiangsu/q17_fig1.png}{8.0cm}{87bp 63bp 73bp 36bp}
\bitmapfigure{img/2014/jiangsu/q17_fig1.png}
\end{problem}

\begin{solution}
记椭圆焦半距为 \(c\)，则 \(c^2=a^2-b^2\)，且 \(F_1=(-c,0)\)，\(F_2=(c,0)\).

\begin{enumerate}
\item 由 \(BF_2=\sqrt2\)，得
\[
b^2+c^2=2.
\]
又因为点 \(C\left(\frac43,\frac13\right)\) 在椭圆上，且 \(a^2=b^2+c^2=2\)，所以
\[
\frac{16}{9\cdot2}+\frac{1}{9b^2}=1
\]
解得 \(b^2=1\). 因而 \(c^2=1\)，\(a^2=2\)，椭圆的方程为
\[
\frac{x^2}{2}+y^2=1.
\]

\item 设 \(C=(x,y)\)，由图形位置可知 \(x>c\)，\(y>0\)，且 \(A=(x,-y)\). 直线 \(BF_2\) 的斜率由两种表示相等，得
\[
\frac{-y-b}{x}=-\frac{b}{c}
\]
即 \(x=c\left(1+\frac{y}{b}\right)\). 令 \(t=\frac{y}{b}>0\)，则
\[
x=c(1+t).
\]
设椭圆的离心率为 \(e=\frac ca\). 由点 \(C\) 在椭圆上，得
\[
e^2(1+t)^2+t^2=1.
\]

由 \(F_1C\perp AB\)，有 \(\overrightarrow{F_1C}=(x+c,y)\)，\(\overrightarrow{AB}=(-x,b+y)\).
所以
\[
y(b+y)=x(x+c).
\]
代入 \(x=c(1+t)\)，\(y=bt\)，并约去正数 \(1+t\)，得
\[
b^2t=c^2(t+2)
\]
所以
\[
t=\frac{2c^2}{b^2-c^2}=\frac{2e^2}{1-2e^2}.
\]
令 \(u=e^2\). 因为 \(t>0\)，所以 \(1-2u>0\)，从而 \(1+t=\frac{1}{1-2u}\). 代入点 \(C\) 的椭圆方程得
\[
\frac{u}{(1-2u)^2}+\frac{4u^2}{(1-2u)^2}=1
\]
即 \(4u^2+u=(1-2u)^2\)，从而 \(u=\frac15\). 因为 \(e>0\)，所以
\[
e=\frac{\sqrt5}{5}.
\]
\end{enumerate}
\end{solution}

\begin{problem}
如图，为了保护河上古桥 \(OA\)，规划建一座新桥 \(BC\)，同时设立一个圆形保护区. 规划要求：新桥 \(BC\) 与河岸 \(AB\) 垂直；保护区的边界为圆心 \(M\) 在线段 \(OA\) 上并与 \(BC\) 相切的圆，且古桥两端 \(O\) 和 \(A\) 到该圆上任意一点的距离均不少于 \(80 \mathrm{~m}\). 经测量，点 \(A\) 位于点 \(O\) 正北方向 \(60 \mathrm{~m}\) 处，点 \(C\) 位于点 \(O\) 正东方向 \(170 \mathrm{~m}\) 处 （\(OC\) 为河岸）， \(\tan \angle BCO = \frac{4}{3}\).

\begin{enumerate}
\item 求新桥 \(BC\) 的长；
\item 当 \(OM\) 多长时，圆形保护区的面积最大？
\end{enumerate}

\bitmapfigure{img/2014/jiangsu/q18_fig1.png}
\end{problem}

\begin{solution}
建立如图所示的直角坐标系，取 \(O=(0,0)\)，则 \(A=(0,60)\)，\(C=(170,0)\). 由 \(\tan\angle BCO=\frac43\)，设存在 \(t>0\)，使得
\[
B=(170-3t,4t).
\]

\begin{enumerate}
\item 由 \(BC\perp AB\)，有
\[
(170-3t,4t-60)\cdot(3t,-4t)=0.
\]
化简得 \(t(750-25t)=0\). 因为 \(t>0\)，所以 \(t=30\)，从而
\[
BC=\sqrt{(3t)^2+(4t)^2}=5t=150 \mathrm{~m}.
\]

\item 设 \(M=(0,m)\)，则 \(0\leqslant m\leqslant60\). 由 \(B=(80,120)\)，直线 \(BC\) 的方程为
\[
4x+3y-680=0.
\]
所以圆的半径为圆心到直线 \(BC\) 的距离
\[
r=\frac{680-3m}{5}=136-\frac35m.
\]

又有 \(r\geqslant100>60\geqslant m\)，\(r\geqslant100>60\geqslant60-m\)，所以圆上各点到 \(O\) 和 \(A\) 的最小距离分别为 \(r-m\) 和 \(r-(60-m)\). 由题设
\(r-m\geqslant80\)，\(r-(60-m)\geqslant80\).
代入 \(r=136-\frac35m\)，得
\[
10\leqslant m\leqslant35.
\]
圆形保护区的面积为 \(\pi r^2\)，而 \(r\) 随 \(m\) 增大而减小，所以面积最大时 \(m\) 取最小值，即
\[
OM=10 \mathrm{~m}.
\]
\end{enumerate}
\end{solution}

\begin{problem}
已知函数 \( f(x) = \e^{x} + \e^{-x} \)，其中 \( \e \) 是自然对数的底数.
\begin{enumerate}
\item 证明： \( f(x) \) 是 \(\R\) 上的偶函数；
\item 若关于 \(x\) 的不等式 \(m f(x) \leqslant \e^{-x} + m - 1\) 在 \((0, +\infty)\) 上恒成立，求实数 \(m\) 的取值范围；
\item 已知正数 \(a\) 满足：存在 \(x_0 \in [1, +\infty)\)，使得 \(f(x_0) < a(-x_0^3 + 3x_0)\) 成立. 试比较 \(\e^{a - 1}\) 与 \(a^{\e - 1}\) 的大小，并证明你的结论.
\end{enumerate}
\end{problem}

\begin{solution}
\begin{enumerate}
\item 对任意 \(x\in\R\)，有
\[
f(-x)=\e^{-x}+\e^x=f(x)
\]
所以 \(f(x)\) 是 \(\R\) 上的偶函数.

\item 令 \(t=\e^x\)，因为 \(x>0\)，所以 \(t>1\). 原不等式等价于
\[
m\left(t+\frac1t\right)\leqslant\frac1t+m-1
\]
即
\[
m\left(t+\frac1t-1\right)\leqslant\frac1t-1.
\]
由于 \(t+\frac1t-1>0\)，所以
\[
m\leqslant\frac{\frac1t-1}{t+\frac1t-1}=-\frac{t-1}{t^2-t+1}.
\]
令 \(s=t-1>0\)，则
\[
-\frac{t-1}{t^2-t+1}=-\frac{s}{s^2+s+1}\geqslant-\frac13
\]
其中不等号等价于 \((s-1)^2\geqslant0\)，且当 \(s=1\) 时取等号. 因而上式对任意 \(t>1\) 成立的充要条件为
\[
m\leqslant-\frac13.
\]

\item 因为 \(f(x_0)>0\)，题设不等式右端必须为正，所以 \(x_0\in[1,\sqrt3)\). 当 \(x\geqslant1\) 时，\(f'(x)=\e^x-\e^{-x}>0\)，故 \(f(x)\) 单调递增，从而
\[
f(x_0)\geqslant f(1)=\e+\e^{-1}.
\]
又因为
\[
3x_0-x_0^3\leqslant2
\]
这是因为 \(x_0^3-3x_0+2=(x_0-1)^2(x_0+2)\geqslant0\)，所以
\[
a>\frac{f(x_0)}{3x_0-x_0^3}\geqslant\frac{\e+\e^{-1}}2>1.
\]

令 \(\Phi(a)=a-1-(\e-1)\ln a\)（\(a>1\)）.
则
\[
\Phi'(a)=1-\frac{\e-1}{a}.
\]
因此 \(\Phi\) 在 \((1,\e-1)\) 上递减，在 \((\e-1,+\infty)\) 上递增；又 \(\Phi(1)=\Phi(\e)=0\). 由此可得，当 \(1<a<\e\) 时 \(\Phi(a)<0\)；当 \(a=\e\) 时 \(\Phi(a)=0\)；当 \(a>\e\) 时 \(\Phi(a)>0\).
而
\[
\ln\frac{\e^{a-1}}{a^{\e-1}}=\Phi(a)
\]
所以
\[
\begin{cases}
\e^{a-1}<a^{\e-1},&1<a<\e,\\
\e^{a-1}=a^{\e-1},&a=\e,\\
\e^{a-1}>a^{\e-1},&a>\e.
\end{cases}
\]
\end{enumerate}
\end{solution}

\begin{problem}
设数列 \(\{a_{n}\}\) 的前 \(n\) 项和为 \(S_{n}\). 若对任意正整数 \(n\)，总存在正整数 \(m\)，使得 \(S_{n} = a_{m}\)，则称 \(\{a_{n}\}\) 是“\(H\) 数列”.
\begin{enumerate}
\item 若数列 \( \{a_{n}\} \) 的前 \(n\) 项和 \( S_{n}=2^{n} \)（\(n\in \N^{*}\)），证明： \( \{a_{n}\} \) 是“\(H\) 数列”；
\item 设\(\{a_{n}\}\) 是等差数列，其首项 \(a_{1} = 1\)，公差 \(d < 0\). 若 \(\{a_{n}\}\) 是“\(H\) 数列”，求 \(d\) 的值；
\item 证明：对任意的等差数列 \(\{a_{n}\}\)，总存在两个“\(H\) 数列” \(\{b_{n}\} \) 和 \(\{c_{n}\}\)，使得 \(a_{n} = b_{n} + c_{n}\)（\(n \in \N^{*}\)）成立.
\end{enumerate}
\end{problem}

\begin{solution}
\begin{enumerate}
\item 当 \(n\geqslant2\) 时，
\[
a_n=S_n-S_{n-1}=2^n-2^{n-1}=2^{n-1}.
\]
并且 \(a_1=S_1=2\). 对任意正整数 \(n\)，有
\[
S_n=2^n=a_{n+1}.
\]
因此 \(\{a_n\}\) 是 \(H\) 数列.

\item 由 \(a_1=1\)，公差为 \(d\)，得 \(a_2=1+d\)，且
\[
S_2=2+d.
\]
由于数列是 \(H\) 数列，存在正整数 \(m\)，使得
\[
2+d=a_m=1+(m-1)d.
\]
所以 \(1=(m-2)d\). 因为 \(d<0\)，故 \(m<2\). 又 \(m\) 为正整数，所以 \(m=1\)，从而 \(d=-1\).

反之，当 \(d=-1\) 时，\(a_n=2-n\)，且
\[
S_n=n-\frac{n(n-1)}2=\frac{n(3-n)}2\leqslant1.
\]
令 \(m=2-S_n\)，则 \(m\) 为正整数，且
\[
a_m=2-m=S_n.
\]
因此 \(d=-1\) 时数列确为 \(H\) 数列，故所求公差为 \(d=-1\).

\item 任取等差数列 \(\{a_n\}\)，设其通项为
\[
a_n=a_1+(n-1)d.
\]
定义 \(b_n=(d-a_1)(n-1)\)，\(c_n=a_1n\).
则
\[
b_n+c_n=(d-a_1)(n-1)+a_1n=a_1+(n-1)d=a_n.
\]
对任意正整数 \(n\)，
\[
\sum_{k=1}^n b_k=(d-a_1)\frac{n(n-1)}2=b_{1+\frac{n(n-1)}2},
\]
其中 \(1+\frac{n(n-1)}2\) 是正整数；同理，
\[
\sum_{k=1}^n c_k=a_1\frac{n(n+1)}2=c_{\frac{n(n+1)}2},
\]
其中 \(\frac{n(n+1)}2\) 是正整数. 因而 \(\{b_n\}\) 和 \(\{c_n\}\) 都是 \(H\) 数列，且 \(a_n=b_n+c_n\).
\end{enumerate}
\end{solution}

\section{附加题}

\begin{problem}[A]
如图，\(AB\) 是圆 \(O\) 的直径，\(C\)，\(D\) 是圆 \(O\) 上位于 \(AB\) 异侧的两点，证明： \( \angle OCB = \angle D \).

\bitmapinlinefigure[0.5\linewidth][]{img/2014/jiangsu/q21_fig1.png}
\end{problem}

\begin{solution}
因为 \(OC=OB\)，所以三角形 \(OCB\) 是等腰三角形，从而
\[
\angle OCB=\angle OBC.
\]
又因为 \(A\)，\(O\)，\(B\) 共线，所以 \(\angle OBC=\angle ABC\). 角 \(ABC\) 与角 \(ADC\) 都是圆周角，且它们所对的弦都是 \(AC\)，故
\[
\angle ABC=\angle ADC=\angle D.
\]
因此 \(\angle OCB=\angle D\).
\end{solution}

\reuseproblemnumber
\begin{problem}[B]
已知矩阵 \(A=\begin{bmatrix}-1&2\\1&x\end{bmatrix}\)，\(B=\begin{bmatrix}1&1\\2&-1\end{bmatrix}\)，向量 \(\alpha=\begin{bmatrix}2\\y\end{bmatrix}\)，\(x\)，\(y\) 为实数，若 \( A\alpha=B\alpha \)，求 \( x+y \) 的值.
\end{problem}

\begin{solution}
由 \(A\alpha=B\alpha\)，分别计算两边得
\[
A\alpha=\begin{bmatrix}-1&2\\1&x\end{bmatrix}\begin{bmatrix}2\\y\end{bmatrix}=\begin{bmatrix}2y-2\\2+xy\end{bmatrix}
\]
以及
\[
B\alpha=\begin{bmatrix}1&1\\2&-1\end{bmatrix}\begin{bmatrix}2\\y\end{bmatrix}=\begin{bmatrix}2+y\\4-y\end{bmatrix}.
\]
比较对应分量，得 \(2y-2=2+y\)，\(2+xy=4-y\).
第一式得 \(y=4\)，代入第二式得 \(2+4x=0\)，所以 \(x=-\frac12\). 因而
\[
x+y=\frac72.
\]
\end{solution}

\reuseproblemnumber
\begin{problem}[C]
在平面直角坐标系 \(xOy\) 中，已知直线 \(l\) 的参数方程为 \(\begin{cases}
x = 1 - \frac{\sqrt{2}}{2} t,\\
y = 2 + \frac{\sqrt{2}}{2} t,
\end{cases}\)（\(t\) 为参数），直线 \(l\) 与抛物线 \(y^2 = 4x\) 交于 \(A\)，\(B\) 两点，求线段 \(AB\) 的长.
\end{problem}

\begin{solution}
由参数方程两式相加，消去参数 \(t\)，得直线 \(l\) 的方程
\[
x+y=3.
\]
与抛物线方程联立，得
\[
(3-x)^2=4x
\]
即
\[
x^2-10x+9=0.
\]
所以 \(x=1\) 或 \(x=9\)，相应地 \(y=2\) 或 \(y=-6\). 两交点为 \((1,2)\)，\((9,-6)\)，故
\[
AB=\sqrt{(9-1)^2+(-6-2)^2}=8\sqrt2.
\]
\end{solution}

\reuseproblemnumber
\begin{problem}[D]
已知 \(x > 0\)，\(y > 0\)，证明： \((1 + x + y^2)(1 + x^2 + y) \geqslant 9xy\).
\end{problem}

\begin{solution}
由基本算术平均不等式，得
\[
1+x+y^2\geqslant3\sqrt[3]{xy^2}
\]
以及
\[
1+x^2+y\geqslant3\sqrt[3]{x^2y}.
\]
两式相乘，得
\[
(1+x+y^2)(1+x^2+y)\geqslant9\sqrt[3]{x^3y^3}=9xy.
\]
\end{solution}

\begin{problem}
盒中共有 \(9\) 个球，其中有 \(4\) 个红球，\(3\) 个黄球和 \(2\) 个绿球，这些球除颜色外完全相同.
\begin{enumerate}
\item 从盒中一次随机取出 \(2\) 个球，求取出的 \(2\) 个球颜色相同的概率 \(P\)；
\item 从盒中一次随机取出 \(4\) 个球，其中红球，黄球，绿球的个数分别记为 \(x_1\)，\(x_2\)，\(x_3\)，随机变量 \(X\) 表示 \(x_1\)，\(x_2\)，\(x_3\) 中的最大数，求 \(X\) 的概率分布和数学期望 \(E(X)\).
\end{enumerate}
\end{problem}

\begin{solution}
从 \(9\) 个球中任取 \(4\) 个球的基本事件总数为
\[
\mathrm C_9^4=126.
\]

\begin{enumerate}
\item 从 \(9\) 个球中任取 \(2\) 个球，取法总数为 \(\mathrm C_9^2\). 颜色相同的取法数为
\[
\mathrm C_4^2+\mathrm C_3^2+\mathrm C_2^2=6+3+1=10.
\]
所以
\[
P=\frac{10}{\mathrm C_9^2}=\frac5{18}.
\]

\item 记取出的 \(4\) 个球中红、黄、绿球的个数分别为 \(x_1\)，\(x_2\)，\(x_3\). 当 \(X=4\) 时，只能取出 \(4\) 个红球，所以对应取法数为
\[
N_4=\mathrm C_4^4=1.
\]

当 \(X=3\) 时，可能的颜色个数组合为 \((3,1,0)\)，\((3,0,1)\)，\((1,3,0)\)，\((0,3,1)\)，故
\[
N_3=\mathrm C_4^3\mathrm C_3^1+\mathrm C_4^3\mathrm C_2^1+\mathrm C_4^1\mathrm C_3^3+\mathrm C_3^3\mathrm C_2^1=12+8+4+2=26.
\]

当 \(X=2\) 时，可能的颜色个数组合为 \((2,2,0)\)，\((1,2,1)\)，\((2,1,1)\)，\((0,2,2)\)，\((1,1,2)\)，\((2,0,2)\).
因此
\[
\begin{aligned}
N_2={}&\mathrm C_4^2\mathrm C_3^2+\mathrm C_4^1\mathrm C_3^2\mathrm C_2^1+\mathrm C_4^2\mathrm C_3^1\mathrm C_2^1\\
&+\mathrm C_3^2\mathrm C_2^2+\mathrm C_4^1\mathrm C_3^1\mathrm C_2^2+\mathrm C_4^2\mathrm C_2^2\\
={}&18+24+36+3+12+6=99.
\end{aligned}
\]
四个球来自三种颜色，所以 \(X\) 不可能取 \(1\). 于是 \(X\) 的概率分布为

\begin{center}
\begin{tabular}{c|cccc}
\(k\) & \(1\) & \(2\) & \(3\) & \(4\) \\
\hline
\(P(X=k)\) & \(0\) & \(\frac{99}{126}=\frac{11}{14}\) & \(\frac{26}{126}=\frac{13}{63}\) & \(\frac1{126}\)
\end{tabular}
\end{center}

所以
\[
E(X)=2\cdot\frac{99}{126}+3\cdot\frac{26}{126}+4\cdot\frac1{126}=\frac{20}{9}.
\]
\end{enumerate}
\end{solution}

\begin{problem}
已知函数 \( f_{0}(x) = \frac{\sin x}{x} \) （\(x > 0\)），设 \(f_{n}(x)\) 为 \(f_{n-1}(x)\) 的导数，\(n \in \N^{*}\).
\begin{enumerate}
\item 求 \(2f_{1}\left(\frac{\pi}{2}\right) + \frac{\pi}{2}f_{2}\left(\frac{\pi}{2}\right)\) 的值；
\item 证明：对任意的 \(n \in \N^{*}\)，等式 \(\left|n f_{n-1}\left(\frac{\pi}{4}\right) + \frac{\pi}{4}f_{n}\left(\frac{\pi}{4}\right)\right| = \frac{\sqrt{2}}{2}\) 都成立.
\end{enumerate}
\end{problem}

\begin{solution}
由 \(xf_0(x)=\sin x\)，对等式两边求一次和两次导数，分别得
\[
f_0(x)+xf_1(x)=\cos x
\]
\[
2f_1(x)+xf_2(x)=-\sin x.
\]
因此
\[
2f_1\left(\frac{\pi}{2}\right)+\frac{\pi}{2}f_2\left(\frac{\pi}{2}\right)=-1.
\]

对等式 \(xf_0(x)=\sin x\) 求 \(n\) 阶导数. 由乘积求导公式，左边为
\[
\frac{\mathrm d^n}{\mathrm d x^n}\left(xf_0(x)\right)=xf_n(x)+nf_{n-1}(x)
\]
右边为
\[
\frac{\mathrm d^n}{\mathrm d x^n}(\sin x)=\sin\left(x+\frac{n\pi}{2}\right).
\]
所以对任意正整数 \(n\)，均有
\[
nf_{n-1}(x)+xf_n(x)=\sin\left(x+\frac{n\pi}{2}\right).
\]
令 \(x=\frac\pi4\)，得
\[
\left|nf_{n-1}\left(\frac\pi4\right)+\frac\pi4f_n\left(\frac\pi4\right)\right|=\left|\sin\left(\frac\pi4+\frac{n\pi}{2}\right)\right|=\frac{\sqrt2}{2}.
\]
\end{solution}
